If $x+3>0$, a real power equals $1$ exactly when the base is $1$ or the exponent is $0$. The first gives $x=-2$. The second gives $(x+5)(x-1)=0$, of which only $x=1$ lies in this positive-base case.

At $x=-3$, the expression is $0^{-8}$ and is undefined. If $x+3<0$ and $|x+3|\ne1$, any defined real power equal to $1$ must have exponent zero, giving $x=-5$, which indeed yields $(-2)^0=1$. Finally $x+3=-1$ means $x=-4$, and the exponent is $-5$, giving $-1$, not $1$.

Therefore the complete solution set is $\boxed{\{-5,-2,1\}}$.
